Before introducing the mutation graph, we fix normalization conditions for the Laurent polynomials that we consider.

\par\noindent\textbf{Definition 2.2.}  \label{def:normalization}
A Laurent polynomial $f\in\mathbb{C}[N]$ is \emph{normalized} if for all vertices $v$ of $\operatorname{Newt}(f)$, the coefficient of the monomial $\bm{x}^v$ in $f$ is~$1$.

\par

\par\noindent\textbf{Convention 2.3.} 
From here onwards, we assume that all Laurent polynomials (and all mutation factors) are normalized. Similarly, although our Laurent polynomials are defined over $\mathbb{C}$, our expectation is that Laurent polynomials that are mirror to Fano manifolds have coefficients that are non-negative integers. From here onwards, we will require that all Laurent polynomials (and all mutation factors) have non-negative integer coefficients. Recall our standing convention that if $f \in \mathbb{C}[N]$ then the exponents of monomials in $f$ generate $N$.

\par
